% Display aliases only: the frozen rows, circuit identities, and numbers are unchanged.
\ExplSyntaxOn
\NewExpandableDocumentCommand{\PaperCircuitLabel}{m}{
  \str_case_e:nnF {#1} {
    {qft_n18}{QFT-18}
    {ising_n42}{Ising-42}
    {ising_model_16}{Ising-16}
    {ising_model_13}{Ising-13}
  } {#1}
}
\ExplSyntaxOff
\begin{table}[!t]
\caption{典型高增益电路的模型保真度比较}
\label{tab:representative-circuits}
\centering
\PaperTableSetup
\let\PaperRawTexttt\texttt
\renewcommand{\texttt}[1]{\PaperRawTexttt{\PaperCircuitLabel{#1}}}
\begin{tabularx}{\columnwidth}{@{}lLcrrr@{}}
\toprule
基准集 & 电路 & $n/g_2$ & ZAC & ICCAD & GA-LK\\
\midrule
\RepresentativeCircuitRows
\bottomrule
\end{tabularx}
\PaperTableNote{注：电路简称与图~\ref{fig:experimental-summary}一致；$n/g_2$为量子比特数/双比特门数。ICCAD指ICCAD/QMAP，粗体为该电路的最高模型保真度。选例条件见正文，总体结果见表~\ref{tab:main-results}。}
\end{table}
